\chapter{Estado de la Cuestión}
\label{cap:estadoDeLaCuestion}

\chapterquote{En primer lugar, nada nace de la nada; pues, si así fuera, cualquier cosa nacería de cualquier otra, sin necesidad de semilla alguna.}{Epicuro}

El motor de exploración descrito en este trabajo se apoya en una tradición de
décadas de investigación, las cuales indagan en cómo permitir que un usuario explore colecciones digitales de tamaño
considerable sin conocer de antemano lo que busca, y se inscribe igualmente en la línea de los modelos
de hipermedia que articularon, mucho antes de la web tal y como la conocemos, las primitivas básicas
para representar información interconectada. Este capítulo recorre ese marco con dos objetivos
complementarios. En primer lugar, identificar qué decisiones técnicas del proyecto son herederas
directas de aportaciones previas, y cuáles otras se separan deliberadamente de ellas. En segundo lugar,
delimitar el espacio dejado por los trabajos relacionados para justificar la existencia de una nueva
propuesta. La sección~\ref{sec:marcoTeorico} desarrolla el marco teórico, distinguiendo la navegación
basada en etiquetas de los modelos de hipermedia más estructurados. La sección~\ref{sec:trabajosRelacionados}
revisa los proyectos con los que el sistema está en diálogo más estrecho, en particular la plataforma Clavy
y propuestas recientes basadas en grandes modelos del lenguaje.

\section{Marco teórico}
\label{sec:marcoTeorico}

La pregunta de fondo que recorre este capítulo es cómo organizar una colección de objetos heterogénea
de manera que un usuario pueda acercarse a ella de forma exploratoria, sin formular una consulta
exacta de antemano y sin depender exclusivamente de un sistema de recomendación opaco. El recorrido
histórico de esa pregunta se ha articulado, a grandes rasgos, en torno a dos paradigmas. El primero, de
inspiración bibliotecaria, organiza los objetos a través de descripciones controladas (etiquetas, atributos,
metadatos) y permite recorrerlos refinando progresivamente esas descripciones. El segundo, nacido en el
ámbito de la documentación interconectada, modela la colección como un grafo de fragmentos enlazados
y cuyo mecanismo primario de exploración se centra en la transición entre nodos. Nuestro motor toma
elementos de ambos. Aplica filtrado por metadatos en el sentido del primer paradigma, y permite saltar
entre colecciones siguiendo relaciones tipadas, en línea con el segundo.

\subsection{Navegación basada en etiquetas de grandes colecciones de objetos}
\label{subsec:tagBasedBrowsing}

La navegación basada en etiquetas, o \textit{tag-based browsing}, parte de un principio sencillo. Cada
objeto de la colección se describe mediante un conjunto de pares atributo-valor, también llamados
etiquetas, y la interacción del usuario consiste en seleccionar una o varias de esas etiquetas para
restringir progresivamente el conjunto de objetos visibles. La selección de una etiqueta no solo filtra
los resultados, sino que reduce también el espacio de etiquetas disponibles, presentando al usuario
únicamente aquellas que pueden seguir refinando el conjunto activo. Este patrón tiene su origen formal
en el trabajo seminal de \citet{godin1986lattice}, que modela el espacio de navegación como un retículo
de conceptos, donde cada estado de exploración corresponde a un par formado por un conjunto de
objetos y un conjunto de etiquetas que los caracteriza, y las transiciones entre estados se obtienen
añadiendo o eliminando etiquetas. Esta intuición resurge con fuerza en los \textit{semantic file systems}
descritos por \citet{gifford1991semantic}, donde el sistema de ficheros abandona la jerarquía rígida de
directorios y la sustituye por una organización plana de etiquetas que se combinan a demanda para
generar vistas consultables como si fueran carpetas virtuales.

Sobre estas dos referencias se construye buena parte de la literatura posterior. La idea común es que el
usuario opera sobre un \textit{estado de navegación} caracterizado por el conjunto de etiquetas activas en
ese momento, al que se asocian dos derivados: el conjunto de objetos filtrados por dichas etiquetas, y el
conjunto de etiquetas seleccionables que pueden añadirse para reducir aún más esos objetos sin
vaciarlos. Cada acción del usuario, ya sea añadir o eliminar una etiqueta, transforma el estado y obliga
a recalcular ambos conjuntos derivados. La formalización de este modelo y su materialización
algorítmica han sido el objeto de un cuerpo notable de trabajo, en el que se sitúa de forma directa
la línea de investigación que dio lugar a la plataforma Clavy \citep{gayoso2016siie}, predecesora más
inmediata de nuestra propuesta. Los autores formalizan la exploración como un autómata de navegación cuyos 
estados son conjuntos de objetos y cuyas transiciones están etiquetadas con pares atributo-valor.

Este modelo tiene una propiedad relevante para nuestro trabajo, y que conviene subrayar antes de seguir.
La caracterización del filtrado como una conjunción acumulativa de condiciones se traduce de forma
natural a una consulta sobre una base de datos relacional o sobre un índice invertido, dado que cada
etiqueta seleccionada corresponde a una condición \texttt{AND} adicional. Nuestra implementación
hereda exactamente esa estructura. Los filtros de metadatos descritos en el capítulo de
desarrollo se construyen como cláusulas \texttt{EXISTS} encadenadas, y el conjunto de etiquetas
seleccionables, que en nuestro sistema denominamos \textit{facetas}, se computa de manera dinámica
agregando los valores presentes en los objetos visibles. La diferencia, sin embargo, es que la literatura
clásica de navegación por etiquetas se ha centrado mayoritariamente en colecciones planas, en las que
todos los objetos son del mismo tipo y comparten un único espacio de etiquetas. En escenarios de
mayor riqueza semántica, donde conviven entidades de naturaleza distinta (películas y personas, libros
y autores, productos y fabricantes), el modelo plano resulta limitado. Allí donde la literatura propone
una única colección descrita por etiquetas heterogéneas, nuestro motor distingue colecciones múltiples
descritas cada una por sus propios atributos y conectadas entre sí por relaciones explícitas, lo que abre
la puerta al segundo paradigma del que hablábamos al inicio de la sección.

\subsection{Modelos de hipermedia y navegación estructurada}
\label{subsec:hypermedia}

La tradición de la navegación por etiquetas describe el qué (qué condiciones se cumplen, qué objetos
las satisfacen) pero deja en segundo plano el cómo (cómo el usuario llega de un objeto al siguiente, qué
relaciones explícitas existen entre ellos). Este vacío lo cubre la tradición de la hipermedia, una línea de
investigación que ha producido una colección de modelos de referencia cuyo objetivo común era
proporcionar un vocabulario formal para describir documentos interconectados. \citet{gruzman2000hypermedia}
ofrecen una revisión panorámica de esta familia de modelos que conviene resumir aquí, porque
varias de sus categorías son directamente trasladables a las decisiones de diseño de nuestro sistema.

El primero y más antiguo de estos modelos es el HAM (\textit{Hypertext Abstract Machine}), concebido
como una máquina servidor de propósito general para almacenar sistemas de hipertexto. Su modelo de
memoria distingue cinco objetos básicos: grafos, contextos, nodos, enlaces entre nodos y atributos
semánticos. La correspondencia entre estas cinco abstracciones y los componentes de nuestro sistema
es notablemente directa, lo que justifica con detalle la afirmación de que el HAM puede leerse como un
antepasado conceptual del diseño aquí presentado. El \textit{grafo} del HAM, entendido como la red
global de información que da soporte al hiperdocumento, encuentra su equivalente en la combinación
de las tablas \texttt{collections} y \texttt{reference} de nuestra base de datos, que en conjunto definen
el universo de las entidades disponibles y de los vínculos tipados que las conectan. Los \textit{nodos},
unidades indivisibles de contenido en el HAM, se corresponden con las filas de la tabla
\texttt{entities}, cada una de las cuales representa una película, una serie, una persona o cualquier otra
unidad atómica del dominio. Los \textit{enlaces} entre nodos son la traducción literal de las filas de la
tabla \texttt{reference}, que registran asociaciones bidireccionales entre entidades acompañadas de una
\textit{razón} (\textit{Director}, \textit{Actor}, etc.) que tipifica la naturaleza del vínculo. Los
\textit{atributos semánticos} del HAM, expresados en aquel modelo como pares clave-valor adheridos
a nodos y enlaces, tienen un reflejo casi idéntico en nuestra tabla \texttt{metadata}, que asocia a cada
entidad un conjunto de pares atributo-valor sobre los que opera el filtrado clásico.

El paralelismo se completa, y resulta especialmente revelador, en el caso del \textit{contexto}.
El contexto del HAM se concibe como un fragmento de datos del grafo que recoge una vista parcial
pertinente para una sesión de usuario, separada del grafo persistente y susceptible de modificarse sin
alterar la información subyacente. Esa es exactamente la función que cumple la clase \texttt{State} de
nuestro sistema. \texttt{State} encapsula los filtros activos, el historial de transiciones y la colección
sobre la que se está operando en un momento dado, viviendo en memoria dentro de la sesión HTTP y
articulando una vista del grafo persistente que es específica de cada usuario. La distinción entre datos
del dominio, persistentes y compartidos, y contexto de exploración, momentáneo y por sesión, una
decisión que se detalla en el capítulo de arquitectura, encuentra aquí su raíz conceptual.

El segundo modelo de referencia, conceptualmente posterior y más elaborado, es el modelo Dexter
(\textit{Dexter Hypertext Reference Model}). Dexter organiza el sistema de hipermedia en tres capas
con interfaces bien definidas. Una capa de almacenamiento, donde residen los componentes
persistentes y sus enlaces. Una capa de ejecución, que coordina las sesiones de usuario y la
instanciación de los componentes activos. Y una capa intracomponente, que describe la estructura
interna de cada componente y se conecta con el resto a través de un mecanismo explícito de
\textit{anclajes} (\textit{anchors}). Esta triple separación encuentra paralelismos en la arquitectura de
nuestro sistema, donde la capa de persistencia (PostgreSQL y la capa DAO), la capa de ejecución
(\texttt{StateManager} y los \textit{servlets} que orquestan la sesión) y la capa de presentación
(\textit{frontend}) cumplen funciones análogas. Más relevante aún es la noción de \textit{enlace} en
Dexter, donde un enlace no es una propiedad implícita de un nodo sino un componente de primera
clase, dotado de identidad propia, posiblemente multidireccional y susceptible de ser él mismo
referenciado. Nuestro objeto \texttt{Link}, que representa una transición entre colecciones e incluye la
dirección, la razón de la relación y una copia del estado en el momento de realizarla, se inscribe en
esta misma tradición. La transición no es un efecto secundario opaco de hacer clic en un elemento,
sino un objeto explícito y persistente, cuya cadena conforma el historial de exploración y se integra
literalmente en la consulta SQL que se construye en cada petición.

Una tercera línea, mucho más reciente y específicamente vinculada al desarrollo web moderno, es la
recuperación de la idea de \textit{Hypermedia As The Engine Of Application State} (HATEOAS),
formulada originalmente por Fielding como uno de los principios constitutivos del estilo arquitectónico
REST. \citet{vepsalainen2024hypermedia} señala que HATEOAS, pese a haber sido el motor conceptual
detrás de REST, es probablemente la parte menos implementada del modelo, y que los sistemas web
mayoritarios (basados en arquitecturas de página única que mantienen el estado en el cliente) han
desplazado el espíritu original. La idea fundamental de HATEOAS es que la representación que el
servidor envía al cliente no contiene únicamente datos, sino también las acciones que el usuario puede
ejecutar a continuación, codificadas como controles hipermedia. El cliente, por tanto, no necesita
conocer de antemano la estructura completa de la API, sino que descubre el espacio de acciones
disponibles a partir de la respuesta a su petición actual. \citet{vepsalainen2024hypermedia} destaca
también la noción de \textit{transclusión}, esto es, la sustitución parcial de fragmentos de una página
por contenido nuevo recibido del servidor, como una de las primitivas características de los
\textit{frameworks} hipermedia recientes (htmx, Datastar, Hotwire), frente a la actualización completa
o al renderizado en cliente.

Examinado bajo esta lente, nuestro sistema satisface el principio en sentido estricto. El
\textit{endpoint} \texttt{GET /api/facets} devuelve, junto a los objetos visibles, el conjunto completo
de acciones legales en ese estado: valores de metadatos seleccionables, entidades disponibles como
filtros de referencia y relaciones tipadas que habilitan saltos a otras colecciones. El cliente no codifica
internamente el modelo de datos ni el espacio de transiciones, sino que presenta lo recibido y traduce
las elecciones del usuario en peticiones a los \textit{endpoints} que el servidor indica, mientras el
contexto de exploración acumulado, motor real del estado de la aplicación, reside íntegramente en el
\texttt{StateManager}. La transclusión, por su parte, se materializa mediante actualizaciones dinámicas
que reemplazan únicamente las regiones afectadas tras cada acción.

Estos dos paradigmas, el de la navegación por etiquetas y el de los modelos hipermedia, se han
mantenido tradicionalmente en compartimentos separados de la literatura. Los sistemas que filtran
colecciones mediante etiquetas raramente describen sus enlaces entre objetos como elementos de
primera clase, y los sistemas hipermedia tienden a primar la navegación entre nodos sobre el
filtrado dentro de un nodo. La síntesis de ambos paradigmas, que es donde se sitúa nuestra
contribución, no está extensamente cubierta en la literatura. Los trabajos que sí han abordado esa
intersección, y de los que nuestro proyecto es heredero, en mayor o menor medida, son los que se revisan
en la siguiente sección.

\section{Trabajos relacionados}
\label{sec:trabajosRelacionados}

Los modelos discutidos en la sección anterior forman el sustrato conceptual sobre el que se asienta
nuestro motor, pero no son trabajos directamente comparables. En esta sección se revisan dos
proyectos cuya distancia con nuestra propuesta es lo suficientemente corta como para que la
comparación sea inevitable. El primero, Clavy, es la línea de investigación que más profundamente
ha explorado la combinación de filtrado por etiquetas y reconfiguración dinámica de esquemas en
colecciones digitales, y constituye nuestro referente inmediato. El segundo, Knowledge Navigator,
representa una corriente reciente que ataca el mismo problema (la exploración de colecciones
extensas) desde una perspectiva radicalmente distinta, basada en grandes modelos del lenguaje. Una
revisión de estos dos polos basta para situar el espacio que nuestra propuesta pretende ocupar.

\subsection{Clavy}
\label{subsec:clavy}

Clavy es una plataforma experimental para la gestión de repositorios de objetos de aprendizaje con
estructuras dinámicamente reconfigurables, desarrollada a lo largo de la última década en la Universidad
Complutense de Madrid \citep{gayoso2016siie}. Su rasgo distintivo, frente a las propuestas estándar
de repositorios de objetos de aprendizaje basadas en esquemas fijos como LOM o IMS CP, es
permitir que cada repositorio defina su propio esquema de metadatos, que ese esquema se construya
de forma inductiva a medida que se cataloga la colección, y que sea reconfigurable en cualquier momento
por los usuarios. Esta filosofía bottom-up acerca a Clavy a las folksonomías y, más en general, a los
sistemas en los que la organización de la información emerge del uso colaborativo y no se impone de
antemano \citep{gayoso2016wise}.

El núcleo conceptual de Clavy, el que más interesa a nuestra discusión, es su modelo de navegación.
Internamente, Clavy abstrae cada objeto del repositorio como un conjunto de pares elemento-valor que
funcionan como etiquetas. Sobre esa abstracción, la plataforma implementa un modelo de
\textit{tag-based browsing} en el sentido descrito en la
sección~\ref{subsec:tagBasedBrowsing}: el usuario añade y elimina etiquetas, el sistema mantiene un
estado caracterizado por las etiquetas activas, y por cada acción se actualizan tanto el conjunto de
objetos filtrados como el conjunto de etiquetas que siguen siendo seleccionables. La originalidad de
Clavy reside en que la jerarquía de elementos del esquema, presentada al usuario como una organización
facetada del espacio de navegación, puede modificarse en cualquier momento sin necesidad de
reorganizar la colección \citep{gayoso2016isd}. Para hacer compatibles ambos requisitos, los autores
demostraron que el comportamiento del sistema podía caracterizarse como un autómata de navegación
finito independiente de la jerarquía. Dado que ese autómata puede, en el peor caso, crecer
exponencialmente con respecto al tamaño de la colección, la línea posterior de la plataforma se ha
dedicado a desarrollar estrategias de indexación y de \textit{caching} que aproximan el comportamiento
deseado sin materializar el autómata por completo. Este conjunto de técnicas incluye una estrategia
basada en la identificación de estados equivalentes (estados que filtran exactamente el mismo conjunto
de objetos, aunque procedan de combinaciones distintas de etiquetas), que permite reducir
sustancialmente el coste de las operaciones de navegación \citep{gayoso2019smartcache}.

En paralelo a la consolidación del modelo de navegación, Clavy ha evolucionado para acomodar
necesidades expresivas más sofisticadas. La incorporación de elementos multivaluados, esto es, de
elementos del esquema que pueden aparecer múltiples veces en la descripción del mismo objeto,
fue un paso importante en escenarios reales como las bibliotecas digitales humanísticas, donde un
objeto típicamente representa una obra con varios contribuidores, cada uno con sus propios atributos
\citep{gayoso2023siie}. La caracterización de estos esquemas como un tipo restringido de gramática
EBNF abre además la posibilidad de extender el modelo a estructuras alternativas y recursivas,
acercándolo al espíritu de los lenguajes generalizados de marcado. La aplicabilidad práctica de la
plataforma se ha demostrado en repositorios reales de humanidades, lingüística y, recientemente,
medicina, donde Clavy ha servido como soporte para generar paquetes de contenido educativo
estandarizado a partir de colecciones médicas existentes \citep{buendia2019medsys}.

La proximidad de Clavy con nuestro proyecto es evidente, y esa proximidad nos obliga a hacer explícitas las
diferencias. Compartimos la idea de un modelo interno
basado en pares atributo-valor con independencia del dominio, la de una navegación guiada por la
selección incremental de filtros, la del cómputo dinámico de las facetas a partir del estado actual y la
de un motor que opera sobre un modelo abstracto desacoplado de la fuente de datos concreta. Las
diferencias, sin embargo, marcan el espacio que nuestra propuesta ocupa. Clavy se centra en colecciones
homogéneas, donde todos los objetos pertenecen al mismo dominio y comparten un único esquema,
con una organización plana que permite reconfigurar la jerarquía de presentación sin afectar a la
estructura subyacente. Nuestro motor, en cambio, asume desde el principio una colección estratificada
en varias subcolecciones de naturaleza distinta (por ejemplo, películas, series, personas, productoras,
cadenas), cada una con su propio espacio de atributos, y vinculadas entre sí por relaciones explícitamente
tipadas. Esa estructura introduce dos primitivas que no tienen análogo directo en Clavy. Por un lado, los
filtros de referencia, que permiten restringir los objetos visibles de una colección por su vinculación con
entidades concretas de otra. Por otro lado, la navegación transversal, que permite trasladarse de una
colección a otra siguiendo una relación y arrastrando el contexto de filtrado acumulado, con un historial
de transiciones que se materializa en subconsultas SQL anidadas. A estas dos primitivas se añade la
operación de unión, que permite al usuario construir varios contextos de exploración de manera
independiente y combinarlos a posteriori para obtener el conjunto de entidades que satisfacen al menos
uno de ellos. Esta operación cubre casos de uso, como la búsqueda de personas que han participado
como directoras en películas de un género o como actrices en series de otro, que no pueden expresarse
con la semántica puramente conjuntiva del filtrado clásico, en la que todas las condiciones acumuladas
deben cumplirse simultáneamente.

Igualmente importante es la diferencia de énfasis. Clavy ha priorizado la reconfigurabilidad dinámica
del esquema y la indexación eficiente, problemas cruciales cuando la colección y su organización
evolucionan en paralelo, y por ello ha invertido un esfuerzo considerable en el desarrollo de
estructuras de índice (autómatas de navegación, dendrogramas, cachés selectivas) que aproximen el
comportamiento ideal del autómata de navegación sin pagar su coste exponencial en el peor caso.
Nuestro proyecto adopta una posición opuesta. Asume un esquema relativamente estable y delega en
PostgreSQL toda la responsabilidad de la indexación, evitando así reimplementar técnicas de
recuperación de información que la base de datos ya resuelve de forma madura. El espacio que se
libera en el lado del motor se invierte en lo que constituye el aporte diferencial de la propuesta, esto
es, la construcción incremental de consultas SQL que reflejen un contexto de navegación enriquecido
con filtros relacionales, transiciones tipadas entre colecciones y combinaciones de conjuntos sobre
contextos independientes. Es en ese desplazamiento del foco, desde el problema de la indexación hacia
el de la traducción dinámica de un estado de exploración complejo a una consulta ejecutable, donde
nuestra propuesta complementa, sin solapar, la línea de trabajo de Clavy.

\subsection{Knowledge Navigator y la exploración asistida por modelos del lenguaje}
\label{subsec:knowledgenav}

En el extremo opuesto del espectro metodológico se sitúan las propuestas recientes que abordan el
mismo problema de fondo, la exploración de colecciones extensas, mediante técnicas basadas en
grandes modelos del lenguaje. Knowledge Navigator, descrito por \citet{katz2024knowledgenav}, es
representativo de esta corriente. Dada una consulta amplia sobre un dominio científico, el sistema
recupera centenares de documentos relevantes, los embebe en un espacio vectorial, los agrupa
mediante clustering probabilístico (modelos de mezcla gaussiana sobre embeddings reducidos con UMAP),
y delega en un modelo de lenguaje las tareas de etiquetar cada \textit{cluster} con un nombre y una
descripción, filtrar aquellos que resultan poco relevantes para la consulta original y organizar los
resultantes en una jerarquía temática de dos niveles que el usuario puede recorrer.

La filiación de Knowledge Navigator con la tradición clásica de \textit{cluster-based browsing} es
declarada por sus propios autores, que sitúan su propuesta como una recuperación moderna del
paradigma Scatter/Gather de los años noventa. Lo que el modelo del lenguaje aporta a esa tradición es,
fundamentalmente, una mejora en la calidad de las representaciones de los \textit{clusters}, durante
mucho tiempo el principal obstáculo para la adopción práctica del paradigma. Los nombres y
descripciones generados por el modelo, frente a las listas de palabras clave extraídas automáticamente
que caracterizaban a los sistemas anteriores, resultan suficientemente legibles como para servir de
mapa al usuario. Las evaluaciones reportadas en el artículo, tanto automáticas como con expertos
humanos, sugieren que el sistema es capaz de cubrir más del setenta por ciento de los subtemas que un
revisor humano habría identificado en una revisión sistemática de literatura, y de generar además
subtemas adicionales no presentes en la revisión.

La distancia con nuestra propuesta es, en este caso, mucho más amplia que en el caso de Clavy, y por
razones casi opuestas. Knowledge Navigator opera sobre colecciones de texto plano (artículos científicos),
sin un modelo relacional explícito, y construye la organización de la colección a partir de propiedades
estadísticas del lenguaje. La organización resultante es probabilística, no necesariamente reproducible
entre ejecuciones, y depende del sesgo del modelo subyacente. Nuestro motor, en cambio, opera sobre
una colección con un modelo de datos explícito (entidades, atributos, relaciones tipadas) y construye la
organización de la exploración mediante operaciones deterministas y trazables sobre ese modelo. La
elección entre ambos enfoques no es de mayor o menor sofisticación, sino de adecuación al problema:
para una colección no estructurada, donde no existe un modelo relacional latente que un sistema
determinista pueda explotar, la aproximación basada en modelos del lenguaje es probablemente la única
viable. Para colecciones donde sí existe ese modelo, como las que motivan nuestro trabajo, prescindir
de él para reducir el problema a uno de \textit{clustering} sobre embeddings supondría desperdiciar
información estructural valiosa y aceptar la opacidad del modelo del lenguaje a cambio de poco.

Esa distinción permite también identificar de manera explícita el hueco en el estado del arte que
justifica la existencia de este proyecto. Por un lado, las herramientas clásicas de navegación basada en
etiquetas, ejemplificadas por la línea Clavy, ofrecen un modelo determinista y trazable, pero se han
mantenido predominantemente en el escenario de colecciones homogéneas, sin primitivas explícitas
para la navegación transversal entre colecciones de distinta naturaleza ni para la combinación de
contextos de exploración construidos por separado. Por otro lado, las herramientas recientes basadas
en modelos del lenguaje son flexibles ante colecciones no estructuradas, pero renuncian al modelo
relacional explícito y, con él, a la posibilidad de operaciones deterministas sobre el espacio de
navegación. Entre ambos extremos queda un espacio escasamente cubierto de forma
satisfactoria, el de un motor de exploración que mantenga el rigor determinista de la primera
tradición, que extienda el modelo plano de etiquetas con relaciones tipadas entre colecciones y con
operaciones de conjuntos sobre contextos independientes, y que exponga estas capacidades a través
de una API web de manera que sean efectivamente utilizables desde una interfaz interactiva. La
propuesta cuyo desarrollo se detalla en los capítulos siguientes pretende, precisamente, ocupar ese
espacio.
